% Limits of Mathematical Expression
\section{Limits of Mathematical Expression}

\subsection{What Cannot Be Expressed}

\subsubsection{Advanced Arithmetic}
The system cannot express:
\begin{itemize}
    \item Subtraction or negative numbers
    \item Multiplication or division
    \item Fractions or rational numbers
    \item Irrational or transcendental numbers
    \item Algebraic equations
\end{itemize}

\subsubsection{Complex Logic}
Cannot represent:
\begin{itemize}
    \item Quantifiers (for all, there exists)
    \item Nested conditionals
    \item Disjunction (OR) explicitly
    \item Negation
    \item Modal logic systematically
\end{itemize}

\subsubsection{Recursive Structures}
Lacks capacity for:
\begin{itemize}
    \item Self-reference
    \item Unbounded recursion
    \item Nested scopes beyond simple qualification
    \item Tree structures
\end{itemize}

\subsection{Why These Limitations?}

\subsubsection{Functional Sufficiency}
Administrative tasks require only:
\begin{itemize}
    \item Entity identification
    \item Quantity enumeration
    \item Simple relationships
    \item Validation/authorization
\end{itemize}

Advanced mathematics unnecessary for these purposes.

\subsubsection{Cognitive Economy}
Simpler systems offer:
\begin{itemize}
    \item Easier learning and transmission
    \item Faster processing and recognition
    \item Lower error rates
    \item Robust degradation (partial damage)
\end{itemize}

\subsubsection{Technological Constraints}
Medium limitations:
\begin{itemize}
    \item Small seal surface area
    \item Carving/impression technology
    \item No paper or extensive writing surfaces
    \item Durability requirements
\end{itemize}

\subsection{Comparison with Contemporary Systems}

\subsubsection{More Mathematically Sophisticated}
Systems with greater mathematical power:
\begin{itemize}
    \item Mesopotamian sexagesimal numerals
    \item Egyptian fractions
    \item Later Chinese place-value notation
\end{itemize}

\subsubsection{Comparable Complexity}
Systems of similar mathematical scope:
\begin{itemize}
    \item Proto-Elamite numerical tablets
    \item Early Egyptian administrative hieroglyphs
    \item Minoan Linear A administrative records
\end{itemize}

\subsection{Theoretical Boundaries}

\subsubsection{Chomsky Hierarchy}
The Indus system is a regular language (Type 3):
\begin{itemize}
    \item Not context-free (Type 2)
    \item Not context-sensitive (Type 1)
    \item Not recursively enumerable (Type 0)
\end{itemize}

\subsubsection{Computational Complexity}
Recognition complexity:
\begin{itemize}
    \item DFA recognition: $O(n)$ time
    \item Polynomial space (fixed finite)
    \item Decidable in linear time
\end{itemize}

This is the simplest class of formal languages, appropriate for administrative notation.

\subsection{Implications}

These limitations are not failures but design features:
\begin{enumerate}
    \item The system accomplishes its administrative purpose
    \item Simplicity ensures reliability and accessibility
    \item Constraints prevent ambiguity
    \item Limited scope enables rapid processing
\end{enumerate}

The Indus script is optimized for its domain, not general-purpose communication.
